% Part: sets-functions-relations
% Chapter: sets
% Section: subsets

\documentclass[../../../include/open-logic-section]{subfiles}

\begin{document}

\olfileid{sfr}{set}{sub}
\olsection{Subhimpunan lan Himpunan Pangkat}

\begin{explain}
Kita bakal kerep mbandhingake himpunan. Salah siji pambandhing kang
cetha yaiku: \emph{kabeh kang ana ing siji himpunan uga ana ing
himpunan liyane}. Kahanan iki cukup wigati, mula kita perlu
ngenalake notasi anyar.
\end{explain}

\begin{defn}[Subhimpunan]
Yen saben !!{element} himpunan $A$ uga !!a{element} saka~$B$,
mula kita nyebut $A$ minangka \emph{subhimpunan} saka~$B$, lan
nulis $A \subseteq B$. Yen $A$ dudu subhimpunan saka~$B$, kita
nulis $A \not\subseteq B$.
Yen $A \subseteq B$ nanging $A \neq B$, kita nulis $A \subsetneq B$
lan nyebut $A$ minangka \emph{subhimpunan sejati} saka $B$.
\end{defn}

\begin{ex}
Saben himpunan iku subhimpunan saka awake dhewe, lan $\emptyset$
iku subhimpunan saka saben himpunan. Himpunan wilangan natural genep
iku subhimpunan saka himpunan wilangan natural.
Uga, $\{ a, b \} \subseteq \{ a, b, c \}$. Nanging
$\{ a, b, e\}$ dudu subhimpunan saka $\{ a, b, c \}$.
\end{ex}

\begin{ex}
Wilangan $2$ iku !!{element} himpunan wilangan wutuh, dene
himpunan wilangan genep iku subhimpunan saka himpunan wilangan
wutuh. Nanging, sawijining himpunan bisa \emph{bebarengan} dadi
!!a{element} lan subhimpunan saka himpunan liyane, tuladhane
$\{0\} \in \{0, \{0\}\}$ lan uga $\{0\} \subseteq \{0,
\{0\}\}$.
\end{ex}

Ekstensionalitas menehi patokan pepadhan kanggo himpunan:
$A = B$ yen lan mung yen saben !!{element} saka~$A$ uga
!!a{element} saka~$B$ lan kosok baline. Definisi "subhimpunan"
nemtokake $A \subseteq B$ persis minangka separo kapisan saka
patokan iki: saben !!{element} saka~$A$ uga !!a{element} saka~$B$.
Mesthi wae definisi iki uga lumaku yen $A$ lan $B$ diijolake:
yaiku, $B \subseteq A$ yen lan mung yen saben !!{element}
saka~$B$ uga !!a{element} saka~$A$. Iki persis perangan
"kosok baline" ing ekstensionalitas. Kanthi tembung liya,
ekstensionalitas nduweni akibat yen himpunan-himpunan iku padha
yen lan mung yen siji lan sijine padha dadi subhimpunan.

\begin{prop}
$A = B$ yen lan mung yen $A \subseteq B$ lan $B \subseteq A$.
\end{prop}

Saiki uga wektu kang trep kanggo ngenalake sawetara notasi
tambahan kang migunani. Nalika nemtokake kapan $A$ iku subhimpunan
saka~$B$, kita nyebut "saben !!{element} saka~$A$ iku \dots,"
lan ngisi "$\dots$" nganggo "!!a{element} saka $B$".
Nanging \emph{wangun} ukara iki kerep banget digunakake, mula
bakal migunani yen kita ngenalake notasi formal kanggo wangun iki.

\begin{defn}\ollabel{forallxina}
$(\forall x \in A)\phi$ minangka cekakan saka $\forall x(x \in A \lif
\phi)$. Semono uga, $(\exists x \in A)\phi$ minangka cekakan saka
$\exists x(x\in A \land \phi)$.
\end{defn}

Kanthi notasi iki, kita bisa nyatakake yen $A \subseteq B$ yen lan
mung yen $(\forall x \in A)x \in B$.

Saiki kita ngrembug jinis himpunan tartamtu: himpunan kabeh
subhimpunan saka sawijining himpunan.

\begin{defn}[Himpunan Pangkat]
Himpunan kang ngemot kabeh subhimpunan saka himpunan~$A$ diarani
\emph{himpunan pangkat saka}~$A$, ditulis $\Pow{A}$.
  \[
    \Pow{A} = \Setabs{B}{B \subseteq A}
  \]
\end{defn}

\begin{ex}
Apa wae kabeh subhimpunan kang bisa digawe saka $\{ a, b, c \}$?
Yaiku: $\emptyset$, $\{a \}$, $\{b\}$, $\{c\}$, $\{a, b\}$,
$\{a, c\}$, $\{b,c\}$, $\{a, b, c\}$.
Himpunan kabeh subhimpunan iki yaiku $\Pow{\{a,b,c\}}$:
\[
\Pow{\{ a, b, c \}} = \{\emptyset, \{a \}, \{b\}, \{c\}, \{a, b\},
\{b, c\}, \{a, c\}, \{a, b, c\}\}
\]
\end{ex}

\begin{prob}
Dhaptarna kabeh subhimpunan saka $\{a, b, c, d\}$.
\end{prob}

\begin{prob}
Tuduhna yen $A$ duwe $n$ !!{element}s, mula $\Pow{A}$ duwe
$2^n$ !!{element}s.
\end{prob}

\end{document}
